
\color{unidarkgreen}{\subsection{Управление выключателем}}
\color{black}

\begin{enumerate}[label=\arabic{section}.\arabic{subsection}.\arabic{subsubsection}.\arabic*,  align=left, labelsep=4pt, leftmargin=0pt, itemindent=57pt]

\item
Функциональный блок управления выключателем предназначен для формирования сигнала на включение силового выключателя. 

\item
В случае присоединения автотрансформатора к шинам ВН и СН через свой выключатель для каждого из них можно задействовать свою логику управления выключателем. Далее приведено описание логики управления выключателем В ВН, логика управления выключателем В СН аналогична.

\item
Команда на включение выключателя \textit{«Включить В1»} формируется:
\begin{itemize}
\item от команды оперативного включения выключателя (сигнал \textit{«РКВ В1(ЛВ)ВН»});
\item при срабатывании АПВ (сигнал \textit{«Ср. АПВ В1(ЛВ)ВН»}); 
\item от внешнего сигнала включения \textit{«Внеш. вкл. В1(ЛВ)ВН»}.
\end{itemize}

\item
Оперативное включение выключателя осуществляется с контролем наличия разрешающего сигнала от функции контроля синхронизма и напряжений (КСН). Максимальное время, в течение которого продолжается ожидание условий оперативного включения, регулируется уставкой \textbf{«Tожид усл. опер. вкл. В1»}. По истечению данной выдержки времени команда оперативного включения сбрасывается. Выдержка времени на возврат \textbf{«Тпрод опер. вкл. В1»} задает время продления сигнала оперативного включения.

\item
При необходимости может быть введен подхват команды на включение от токового реле, установленного в цепи ЭМВ, в этом случае сигнал на включение удерживается в сработанном состоянии до исчезновения тока через ЭМВ (пока блок-контакты выключателя не разорвут цепь включения). Ввод осуществляется программным ключом \textbf{«Подхват вкл. В1 от РТ ЭМВ»}. 

\item
Предусмотрена возможность блокировки многократных включений на КЗ. Если при наличии сигнала включения будет протекать ток по электромагнитам отключения (что соответствует включению на КЗ), то сигнал \textit{«Включить В1»} заблокируется, пока не пропадут сигналы включения (от АПВ или оперативного включения). Блокировка вводится программным ключом \textbf{«Блок. от многокр. вкл. на КЗ В1»}.

\item
При наличии сигнала \textit{«Блок. включения В1»}, формируемого функциональным блоком контроля силового выключателя, включение выключателя блокируется.

\item
Алгоритм работы функции управления выключателем стороны ВН приведен на рисунке \ref{pic:OPQ}, алгоритм работы функции управления выключателем стороны СН аналогичный. Уставки УВ ВН приведены в таблице \ref{app_set:uvvn}. Уставки УВ СН приведены в таблице \ref{app_set:uvsn}.

\begin{figure}[H]
\centering
\includegraphics[width=1\textwidth,height=1\textheight,keepaspectratio]{\fbpath/761.3391 Управление В/img/OP_Q.pdf}
\captionof{figure}{Функциональная схема логики управления выключателем стороны ВН}\label{pic:OPQ}
\end{figure}


\end{enumerate}